% Volume IX: Integration, Capstones, and Reference Material
% Continuity note for Chapter 54.
% Previous dependencies: This appendix assumes the notation, definitions, and examples developed throughout Chapters 1--53.
% New durable concepts: compact reference tables for notation, matrix derivatives, differential notation, distributions, stochastic processes, information identities, optimization identities, SDE notation, theorem selection, common derivations, and computational recipes.
% Forward hooks: This is the terminal reference chapter for the course.
% Notation commitments: Gradients of scalar losses have the same shape as the variable; Jacobians use numerator layout J_{ij}=partial y_i / partial x_j so dy=J dx; matrix gradients use dL=<grad_W L,dW>_F; stochastic processes are indexed by time and SDEs use dW_t only for Brownian increments.
\section{Chapter 54: Mathematical Appendices}
\label{ch:mathematical-appendices}

This appendix is a working reference, not a replacement for the earlier chapters. Its job is to collect formulas, notation, assumptions, and computational recipes in one place so that the reader can check shapes, identities, and modeling choices while solving problems in computational neuroscience and AI.

\paragraph{Chapter problem.}
Which formulas and procedures should a beginner be able to locate quickly when moving among linear algebra, calculus, probability, information theory, optimization, stochastic processes, machine learning, and neural data analysis?

\paragraph{Section map.}
Section 54.1 is a notation dictionary, with special attention to differential notation and SDE symbols. Section 54.2 gives theorem and identity cards, including distributions, processes, information identities, and optimization identities. Section 54.3 works through common derivations, especially matrix derivatives, softmax cross-entropy, least squares, and shape discipline for backpropagation. Section 54.4 gives computational recipes for common analysis workflows.

\subsection{54.1 Notation dictionary}

\paragraph{Guiding question.}
When the same symbol style appears in many chapters, how do we know what kind of object it denotes and which operations are legal?

\paragraph{Beginner-facing intuition.}
Most mathematical mistakes in this course are object mistakes before they are algebra mistakes. A vector is not a scalar, a density is not a probability, a Jacobian is not a gradient, and \(dW_t\) in an SDE is not an ordinary differential. This dictionary is a compact type-checking guide.

\paragraph{Formal setup and dictionary.}
\begin{center}
\small
\setlength{\tabcolsep}{2pt}
\renewcommand{\arraystretch}{1.16}
\begin{tabular}{p{0.17\linewidth}p{0.21\linewidth}p{0.27\linewidth}p{0.25\linewidth}}
\toprule
Object & Typical notation & Type or shape & Meaning \\
\midrule
Scalar & \(a,\lambda,t,L\) & element of \(\mathbb R\) or subset & one number, time, rate, loss, or parameter \\
Vector & \(x,r,\theta\) & \(\mathbb R^d\) & features, neural response, parameter vector \\
Matrix & \(A,W,X\) & \(\mathbb R^{m\times n}\) & linear map, weights, data matrix \\
Tensor & \(R_{b,t,n}\) & multi-axis array & trial-time-neuron or batch-token-feature data \\
Random variable & \(X,Y,R,S\) & function on \(\Omega\) & uncertain quantity extracted from outcomes \\
Distribution & \(p(x),P_X,q_\theta(y\mid x)\) & probability law & probabilities or densities over values \\
Expectation & \(\mathbb E[X]\) & scalar/vector/matrix shaped like \(X\) & probability-weighted average \\
Gradient & \(\nabla_x L\) & same shape as \(x\) & derivative of scalar \(L\) with respect to \(x\) \\
Jacobian & \(J_f(x)\) & \(\mathbb R^{m\times n}\) for \(f:\mathbb R^n\to\mathbb R^m\) & local linear map \(df_x(h)=J_f(x)h\) \\
Hessian & \(\nabla_x^2 L\) & \(\mathbb R^{n\times n}\) & derivative of a gradient; local curvature \\
\bottomrule
\end{tabular}
\end{center}

\paragraph{Differential notation reference.}
\begin{center}
\small
\setlength{\tabcolsep}{2pt}
\renewcommand{\arraystretch}{1.16}
\begin{tabular}{p{0.17\linewidth}p{0.31\linewidth}p{0.42\linewidth}}
\toprule
Notation & Object & Safe interpretation \\
\midrule
\(dx\) & input perturbation or integration symbol & a formal small change; its meaning depends on context \\
\(df_a\) & linear map & differential of \(f\) at \(a\), mapping perturbations to first-order output changes \\
\(dy/dx\) & scalar derivative & rate of change for \(y=f(x)\), \(x,y\in\mathbb R\) \\
\(\partial f/\partial x_i\) & partial derivative & change in scalar \(f\) as coordinate \(x_i\) varies and others are held fixed \\
\(\nabla f(x)\) & gradient vector & vector satisfying \(df_x(h)=\nabla f(x)^\top h\) for scalar \(f\) \\
\(J_f(x)\) & Jacobian matrix & matrix satisfying \(df_x(h)=J_f(x)h\) for vector-valued \(f\) \\
\(\nabla^2 f(x)\) & Hessian matrix & second-order curvature for scalar \(f\) \\
\(dx/dt=f(x,t)\) & ODE & ordinary time derivative equals deterministic vector field \\
\(dX_t=f(X_t,t)dt+g(X_t,t)dW_t\) & SDE & stochastic update with Brownian increment \(dW_t\), whose variance scales like \(dt\) \\
\bottomrule
\end{tabular}
\end{center}

\paragraph{Core mechanism: local linear type-checking.}
If \(f:\mathbb R^n\to\mathbb R^m\), then for small \(h\in\mathbb R^n\),
\[
f(x+h)=f(x)+J_f(x)h+\text{higher-order error}.
\]
The shapes enforce the meaning:
\[
J_f(x)\in\mathbb R^{m\times n},\qquad h\in\mathbb R^n,\qquad J_f(x)h\in\mathbb R^m.
\]
For a scalar loss \(L:\mathbb R^n\to\mathbb R\), the Jacobian is a row in this convention, but the gradient is stored as a vector \(\nabla L(x)\in\mathbb R^n\) satisfying
\[
dL_x(h)=\nabla L(x)^\top h.
\]

\paragraph{Concrete example.}
Let
\[
f(x_1,x_2)=
\begin{bmatrix}
x_1+x_2\\
x_1x_2\\
x_2^2
\end{bmatrix}.
\]
Then \(f:\mathbb R^2\to\mathbb R^3\), and
\[
J_f(x)=
\begin{bmatrix}
1&1\\
x_2&x_1\\
0&2x_2
\end{bmatrix}.
\]
For \(x=(2,3)\) and \(h=(0.1,-0.2)\),
\[
J_f(x)h=
\begin{bmatrix}
-0.1\\
3(0.1)+2(-0.2)\\
6(-0.2)
\end{bmatrix}
=
\begin{bmatrix}
-0.1\\
-0.1\\
-1.2
\end{bmatrix}.
\]

\paragraph{Computational neuroscience example.}
If \(r=f(s)\in\mathbb R^N\) is a population tuning map from a scalar stimulus \(s\) to \(N\) neural responses, then \(df_s\) is an \(N\times1\) Jacobian. Its entries are tuning-curve slopes. Large slopes relative to noise support precise local stimulus discrimination.

\paragraph{AI or machine-learning example.}
In backpropagation, \(dx\), \(dW\), and \(db\) are perturbations of different shaped objects. A gradient with respect to a weight matrix \(W\in\mathbb R^{m\times n}\) must also have shape \(m\times n\). Shape mismatches often reveal an incorrect transpose or missing batch sum.

\paragraph{Common misconceptions and failure modes.}
Do not treat every \(d\) symbol as an ordinary tiny number. In integrals, \(dx\) marks the variable of integration. In differentials, \(df_x\) is a local linear map. In SDEs, \(dW_t\) is Brownian noise with scale \(\sqrt{dt}\). Also, \(\partial f/\partial x_i\) is not the same object as the full gradient unless all coordinates are collected.

\paragraph{Exercises.}
\begin{enumerate}[label=\textbf{54.1.\arabic*.}, leftmargin=*]
\item \textbf{Mechanical.} If \(f:\mathbb R^5\to\mathbb R^2\), state the shape of \(J_f(x)\). If \(L:\mathbb R^{4\times3}\to\mathbb R\), state the shape of \(\nabla_W L\).
\item \textbf{Proof or derivation.} For scalar \(f:\mathbb R^n\to\mathbb R\), show that the differential formula \(df_x(h)=\nabla f(x)^\top h\) makes \(df_x\) a linear map in \(h\).
\item \textbf{Numerical/computational.} Compute the Jacobian of \(f(x_1,x_2)=(x_1^2,x_1+x_2)\) and evaluate \(J_f(1,2)(3,4)^\top\).
\item \textbf{Interpretation.} In an SDE, why does Brownian noise use \(dW_t\) rather than an ordinary derivative \(dW_t/dt\)?
\item \textbf{Misconception/debugging.} A student says, ``The Hessian and Jacobian are both derivative matrices, so they are interchangeable.'' Correct this using domains, codomains, and shapes.
\end{enumerate}

\subsection{54.2 Core theorem reference}

\paragraph{Guiding question.}
Which theorem or identity should be used for a modeling step, and what assumptions must be checked before using it?

\paragraph{Beginner-facing intuition.}
Theorems are tools with input requirements. Bayes' theorem needs a joint distribution and nonzero evidence. PCA optimality needs squared reconstruction error and linear subspaces. Convex first-order conditions need convexity. Data processing needs a Markov chain. This section lists common tools together with their assumptions and failure modes.

\paragraph{Core theorem and identity cards.}
\begin{center}
\small
\setlength{\tabcolsep}{2pt}
\renewcommand{\arraystretch}{1.12}
\begin{tabular}{p{0.18\linewidth}p{0.30\linewidth}p{0.22\linewidth}p{0.20\linewidth}}
\toprule
Tool & Statement & Use when & Check \\
\midrule
Bayes & \(p(\theta\mid x)=p(x\mid\theta)p(\theta)/p(x)\) & infer latent causes or parameters & evidence nonzero; model specified \\
Least squares & \(X^\top X\widehat w=X^\top y\) & linear regression with squared loss & rank, conditioning, leakage \\
SVD/PCA & top eigenvectors minimize linear reconstruction error & dimensionality reduction & centered data; variance is target \\
Chain rule & \(J_{f\circ g}=J_f(g)J_g\) & backpropagation and sensitivity & shapes and differentiability \\
Gibbs inequality & \(D_{\mathrm{KL}}(p\|q)\geq0\) & distribution mismatch & support of \(q\) covers \(p\) \\
Data processing & \(X\to Y\to Z\Rightarrow I(X;Z)\leq I(X;Y)\) & representation compression & Markov condition \\
Convex optimum & \(\nabla f(x^\star)=0\) implies global optimum for differentiable convex \(f\) & convex optimization & convexity and domain interior \\
Linear stability & eigenvalues of \(J\) have negative real parts & local ODE fixed-point stability & linearization valid locally \\
\bottomrule
\end{tabular}
\end{center}

\paragraph{Common distributions and processes.}
\begin{center}
\small
\setlength{\tabcolsep}{2pt}
\renewcommand{\arraystretch}{1.12}
\begin{tabular}{p{0.17\linewidth}p{0.29\linewidth}p{0.22\linewidth}p{0.22\linewidth}}
\toprule
Object & Formula or defining property & Mean or key parameter & Typical use \\
\midrule
Bernoulli & \(P(X=1)=p\) & \(p\) & binary event or spike/no spike bin \\
Categorical & \(P(X=k)=\pi_k\) & probability vector \(\pi\) & class label or discrete state \\
Poisson & \(P(N=k)=e^{-\lambda}\lambda^k/(k\mathord{!})\) & \(\lambda\) & counts in a window \\
Gaussian & \(p(x)\propto \exp[-(x-\mu)^2/(2\sigma^2)]\) & \(\mu,\sigma^2\) & additive noise, residuals \\
Markov chain & \(P(X_{t+1}=j\mid X_t=i)=P_{ij}\) & transition matrix \(P\) & discrete latent states \\
Brownian motion & independent increments \(\Delta W\sim\mathcal N(0,\Delta t)\) & variance grows as \(t\) & diffusion limit, SDE noise \\
OU process & \(dX_t=-\alpha(X_t-\mu)dt+\sigma dW_t\) & stationary variance \(\sigma^2/(2\alpha)\) & mean-reverting neural noise \\
Poisson process & \(N(t)-N(s)\sim\operatorname{Pois}(\lambda(t-s))\) & rate \(\lambda\) & event times and spike trains \\
\bottomrule
\end{tabular}
\end{center}

\paragraph{Information and optimization identities.}
\begin{center}
\small
\setlength{\tabcolsep}{2pt}
\renewcommand{\arraystretch}{1.12}
\begin{tabular}{p{0.24\linewidth}p{0.36\linewidth}p{0.30\linewidth}}
\toprule
Identity & Formula & Use \\
\midrule
Entropy & \(H(X)=-\sum_x p(x)\log p(x)\) & uncertainty of a discrete variable \\
Cross-entropy & \(H(p,q)=-\sum_x p(x)\log q(x)\) & prediction or coding cost under \(q\) \\
KL divergence & \(D_{\mathrm{KL}}(p\|q)=\sum_xp(x)\log(p(x)/q(x))\) & directed mismatch \\
Mutual information & \(I(X;Y)=D_{\mathrm{KL}}(p(x,y)\|p(x)p(y))\) & dependence or uncertainty reduction \\
First-order optimum & \(\nabla f(x^\star)=0\) & smooth unconstrained optimum candidate \\
Gradient descent & \(x_{k+1}=x_k-\eta\nabla f(x_k)\) & iterative minimization \\
Lagrangian & \(\mathcal L(x,\lambda)=f(x)+\lambda^\top g(x)\) & equality-constrained optimization \\
ELBO & \(\log p(x)\geq \mathbb E_q[\log p(x,z)-\log q(z)]\) & variational inference \\
\bottomrule
\end{tabular}
\end{center}

\paragraph{Core mechanism: theorem selection by assumptions.}
Before applying a theorem, ask four questions:
\[
\text{What are the objects?}\quad
\text{What assumptions are required?}\quad
\text{What conclusion follows?}\quad
\text{What can fail?}
\]
For example, data processing gives \(I(S;Z)\leq I(S;R)\) only when \(S\to R\to Z\) is a Markov chain, such as when \(Z\) is computed from \(R\) without seeing \(S\) directly. If \(Z\) is built using labels \(S\), the assumption may fail.

\paragraph{Concrete example.}
Let \(S\) be a stimulus, \(R\) a neural response, and \(Z=AR\) a linear projection. Since \(Z\) is a deterministic function of \(R\), \(S\to R\to Z\). Therefore \(I(S;Z)\leq I(S;R)\). If a projection is selected using all labels and then evaluated on the same data, the mathematical inequality still holds for the variables, but the empirical estimate of information may be biased upward by selection.

\paragraph{Computational neuroscience example.}
Use a Poisson process theorem for spike times only after checking whether independent increments and conditionally specified intensity are plausible. If refractory periods are strong, a simple homogeneous Poisson process will fail posterior predictive checks even if its mean firing rate is correct.

\paragraph{AI or machine-learning example.}
Use convex optimality for logistic regression with convex loss, but not for a deep network with multiple nonlinear layers. For deep networks, a zero gradient is a stationary point, not a guarantee of global optimality.

\paragraph{Common misconceptions and failure modes.}
Do not use theorem names as decorations. A formula used without its assumptions can be misleading. A lower AIC, higher mutual information estimate, or lower training loss is not automatically a stronger scientific claim. Continuous and discrete formulas often look similar but differ in important ways, especially entropy and densities.

\paragraph{Exercises.}
\begin{enumerate}[label=\textbf{54.2.\arabic*.}, leftmargin=*]
\item \textbf{Mechanical.} For each object, name a suitable distribution or process: binary spike in one bin, spike count in a window, mean-reverting voltage fluctuation, and class label.
\item \textbf{Proof or derivation.} Starting from \(I(X;Y)=H(X)-H(X\mid Y)\), derive the symmetric identity \(I(X;Y)=H(Y)-H(Y\mid X)\) for discrete variables.
\item \textbf{Numerical/computational.} Compute \(D_{\mathrm{KL}}(p\|q)\) in nats for \(p=(1/2,1/2)\) and \(q=(3/4,1/4)\).
\item \textbf{Interpretation.} Explain why a Markov chain assumption is needed before applying data processing.
\item \textbf{Misconception/debugging.} A student uses \(\nabla f(x^\star)=0\) to claim a deep network reached the global optimum. Identify the missing assumption.
\end{enumerate}

\subsection{54.3 Common derivations}

\paragraph{Guiding question.}
Which algebraic derivations occur often enough that we should be able to reproduce and check them quickly?

\paragraph{Beginner-facing intuition.}
Many later formulas are built from a small set of derivative moves: expand a squared norm, use a differential, move factors inside a trace, identify the gradient through the Frobenius pairing, and propagate residuals backward through a chain. This section collects those moves.

\paragraph{Matrix derivative reference.}
For scalar \(L\), matrix gradients are defined by
\[
dL=\langle \nabla_W L,dW\rangle_F
=\operatorname{tr}((\nabla_W L)^\top dW).
\]

\begin{center}
\small
\setlength{\tabcolsep}{2pt}
\renewcommand{\arraystretch}{1.15}
\begin{tabular}{p{0.29\linewidth}p{0.28\linewidth}p{0.33\linewidth}}
\toprule
Function & Gradient & Shape notes \\
\midrule
\(\frac12\|Xw-y\|_2^2\) & \(X^\top(Xw-y)\) & \(X\in\mathbb R^{M\times d}, w\in\mathbb R^d\) \\
\(\frac12\|XW-Y\|_F^2\) & \(X^\top(XW-Y)\) & gradient has shape \(d\times k\) \\
\(\frac12 x^\top A x\), \(A=A^\top\) & \(Ax\) & if \(A\) not symmetric, use \((A+A^\top)x/2\) \\
\(a^\top x\) & \(a\) & linear form gradient \\
\(b^\top W a\) & \(ba^\top\) & \(W\in\mathbb R^{m\times n}\), \(b\in\mathbb R^m\), \(a\in\mathbb R^n\) \\
\(\operatorname{tr}(A^\top W)\) & \(A\) & trace-Frobenius pairing \\
\(-\sum_i y_i\log p_i\), \(p=\operatorname{softmax}(z)\) & \(p-y\) with respect to \(z\) & \(y\) one-hot or probability vector \\
\bottomrule
\end{tabular}
\end{center}

\paragraph{Trace tricks and shape discipline.}
Useful identities include
\[
\langle A,B\rangle_F=\operatorname{tr}(A^\top B),
\]
\[
\operatorname{tr}(ABC)=\operatorname{tr}(BCA)=\operatorname{tr}(CAB)
\]
when shapes are compatible, and
\[
d(Wa)=dW\,a,\qquad d(AW)=A\,dW.
\]
The gradient must have the same shape as the variable. This single rule catches many errors.

\paragraph{Core derivation 1: least squares.}
Let
\[
L(W)=\frac12\|XW-Y\|_F^2,
\qquad E=XW-Y.
\]
Then
\[
dL=\langle E,dE\rangle_F=\langle E,X\,dW\rangle_F.
\]
Move \(X\) to the other side of the Frobenius product:
\[
\langle E,X\,dW\rangle_F
=\operatorname{tr}(E^\top X\,dW)
=\operatorname{tr}((X^\top E)^\top dW)
=\langle X^\top E,dW\rangle_F.
\]
Therefore
\[
\boxed{\nabla_W L=X^\top(XW-Y).}
\]

\paragraph{Core derivation 2: softmax cross-entropy.}
Let
\[
p_i=\frac{e^{z_i}}{\sum_j e^{z_j}},
\qquad
L(z)=-\sum_i y_i\log p_i,
\]
where \(\sum_i y_i=1\). Since
\[
\log p_i=z_i-\log\sum_j e^{z_j},
\]
\[
L(z)=-\sum_i y_i z_i+\log\sum_j e^{z_j}.
\]
Differentiate:
\[
\frac{\partial L}{\partial z_k}
=-y_k+\frac{e^{z_k}}{\sum_j e^{z_j}}
=p_k-y_k.
\]
Thus
\[
\boxed{\nabla_z L=p-y.}
\]

\paragraph{Concrete example.}
For a three-class classifier with
\[
p=(0.7,0.2,0.1),\qquad y=(1,0,0),
\]
the logit gradient is
\[
p-y=(-0.3,0.2,0.1).
\]
The correct class logit is pushed up by gradient descent, because subtracting the negative gradient increases it; the incorrect logits are pushed down.

\paragraph{Computational neuroscience example.}
In a linear neural decoder \(s\approx Rw\), the gradient \(R^\top(Rw-s)\) tells how each neuron's activity correlates with the decoding residual. Large residual correlations indicate weights that should change.

\paragraph{AI or machine-learning example.}
Softmax cross-entropy is the standard classification loss for neural networks. The compact gradient \(p-y\) is why the final-layer backpropagation signal has a simple interpretation: predicted probability minus target probability.

\paragraph{Common misconceptions and failure modes.}
Trace identities cannot ignore shape compatibility. Gradients with respect to matrices are not automatically transposes of familiar vector formulas; derive them with the Frobenius pairing. Softmax and cross-entropy simplify together; differentiating softmax alone is more complicated. Batch averaging changes gradients by a factor of \(1/M\), which must match the implementation.

\paragraph{Exercises.}
\begin{enumerate}[label=\textbf{54.3.\arabic*.}, leftmargin=*]
\item \textbf{Mechanical.} State the gradient of \(\frac12\|Xw-y\|_2^2\) and the shapes of all terms.
\item \textbf{Proof or derivation.} Derive the gradient of \(\frac12 x^\top A x\) when \(A\) is not assumed symmetric.
\item \textbf{Numerical/computational.} For \(p=(0.2,0.5,0.3)\) and \(y=(0,1,0)\), compute the softmax cross-entropy logit gradient.
\item \textbf{Interpretation.} In a neural decoder, what does a positive component of \(R^\top(Rw-s)\) say about a feature's correlation with residual error?
\item \textbf{Misconception/debugging.} A student computes \(\nabla_W \frac12\|XW-Y\|_F^2=(XW-Y)X^\top\). Use shapes to show why this is wrong when \(X\in\mathbb R^{M\times d}\) and \(Y\in\mathbb R^{M\times k}\).
\end{enumerate}

\subsection{54.4 Computational recipes}

\paragraph{Guiding question.}
How can the course's formulas be turned into reliable analysis workflows?

\paragraph{Beginner-facing intuition.}
A computational recipe is a sequence of mathematical decisions: define inputs and outputs, split data, fit parameters, check assumptions, evaluate on held-out data, and report uncertainty. Recipes are not substitutes for understanding, but they prevent common analysis failures.

\paragraph{Recipe cards.}
\begin{center}
\small
\setlength{\tabcolsep}{2pt}
\renewcommand{\arraystretch}{1.12}
\begin{tabular}{p{0.19\linewidth}p{0.35\linewidth}p{0.34\linewidth}}
\toprule
Task & Core steps & Checks \\
\midrule
Train/test model comparison & define observation unit; split; fit only on training; evaluate held-out likelihood or loss & leakage, grouped data, calibration, confidence intervals \\
PCA on neural data & center \(X\); compute covariance or SVD; project; report variance explained; validate behavior readout & scaling, artifacts, held-out reconstruction, overinterpreted plots \\
Poisson GLM for spikes & build design matrix; choose link \(\lambda=\exp(X\beta)\); maximize log-likelihood; test on held-out spikes & refractory effects, overdispersion, time-rescaling \\
Simulate ODE & choose \(x_0,\Delta t\); update \(x_{k+1}=x_k+\Delta t f(x_k,t_k)\); test step size & stability, units, conservation laws \\
Simulate SDE & choose \(X_0,\Delta t\); update \(X_{k+1}=X_k+f\Delta t+g\sqrt{\Delta t}\xi_k\) & noise scaling, random seed, distribution checks \\
Estimate mutual information & define variables and bins or model; estimate entropies; bias-correct or validate by resampling & finite-sample bias, bin choice, continuous variables \\
Posterior predictive check & fit model; simulate replicated data; compare targeted statistics & choose meaningful statistics, not only visual fit \\
\bottomrule
\end{tabular}
\end{center}

\paragraph{Formal setup.}
Most recipes can be written as
\[
\text{data }D
\quad\longrightarrow\quad
\widehat\theta=\operatorname{Fit}(D_{\mathrm{train}})
\quad\longrightarrow\quad
\operatorname{Evaluate}(\widehat\theta,D_{\mathrm{test}})
\quad\longrightarrow\quad
\operatorname{Check}(D,D_{\mathrm{rep}},\widehat\theta).
\]
For stochastic simulations, the recipe is a transition kernel:
\[
X_{k+1}\sim K_{\Delta t}(\cdot\mid X_k).
\]
For Euler-Maruyama applied to
\[
dX_t=f(X_t,t)dt+g(X_t,t)dW_t,
\]
the transition is approximately
\[
X_{k+1}=X_k+f(X_k,t_k)\Delta t+g(X_k,t_k)\sqrt{\Delta t}\xi_k,
\qquad \xi_k\sim\mathcal N(0,I).
\]

\paragraph{Core mechanism: split, fit, simulate, check.}
The same four-step structure appears throughout the course.
\[
\begin{array}{ll}
\text{Split} & \text{protects evaluation from training information.}\\
\text{Fit} & \text{chooses parameters or latent variables under a stated objective.}\\
\text{Simulate or predict} & \text{turns the model into observable consequences.}\\
\text{Check} & \text{compares those consequences with held-out data or diagnostics.}
\end{array}
\]
When a model fails, the failure is useful information: it identifies a missing covariate, wrong noise model, bad preprocessing step, or unsupported scientific claim.

\paragraph{Concrete example.}
To compare two neural encoding models, split stimuli into five folds. For each fold, fit both models on four folds, choose hyperparameters using only training folds or inner validation, and compute held-out log-likelihood on the remaining fold. Average held-out log-likelihoods across folds. Then simulate responses from each fitted model and compare response variance, tuning shape, and residual autocorrelation with real data.

\paragraph{Computational neuroscience example.}
For a spike-train GLM, a recipe is: bin spikes, build stimulus and spike-history covariates, split by trials, fit the Poisson log-likelihood with regularization, evaluate held-out log-likelihood, perform time-rescaling or residual checks, and report uncertainty in filters. Each step has a mathematical object and a failure mode.

\paragraph{AI or machine-learning example.}
For a neural-network classifier, a recipe is: define input and target tensors, split by the scientific generalization target, standardize using training data, train by minibatch gradient descent, monitor validation loss, evaluate held-out calibration and robustness, and inspect errors. The same recipe logic applies to models of biological data and ordinary ML benchmarks.

\paragraph{Common misconceptions and failure modes.}
Do not tune hyperparameters on the test set. Do not report only the best split. Do not compare models with different preprocessing leakage. Do not interpret a simulation match on one statistic as complete model adequacy. Do not simulate an SDE with noise proportional to \(\Delta t\) instead of \(\sqrt{\Delta t}\).

\paragraph{Exercises.}
\begin{enumerate}[label=\textbf{54.4.\arabic*.}, leftmargin=*]
\item \textbf{Mechanical.} Write the Euler-Maruyama update for \(dX_t=-X_tdt+2dW_t\) with step size \(\Delta t\).
\item \textbf{Proof or derivation.} Show why the Brownian increment term in Euler-Maruyama has standard deviation \(\sqrt{\Delta t}\) if \(W_{t+\Delta t}-W_t\sim\mathcal N(0,\Delta t)\).
\item \textbf{Numerical/computational.} Starting from \(X_0=1\), simulate one Euler-Maruyama step for \(dX_t=-X_tdt+dW_t\), \(\Delta t=0.04\), and \(\xi=0.5\).
\item \textbf{Interpretation.} In a posterior predictive check for a spike model, why might matching the mean count be insufficient?
\item \textbf{Misconception/debugging.} A student chooses the model with the lowest test error after trying twenty preprocessing pipelines on the same test set. Explain the problem and give a corrected workflow.
\end{enumerate}

\paragraph{Closing bridge.}
The appendices close the course by turning repeated mathematical moves into a portable reference. The larger lesson is the same one introduced in Chapter 1: name the objects, state the assumptions, check the shapes, derive the mechanism, evaluate on the right data, and keep the scientific claim tied to the mathematical evidence.

\clearpage
